% Vietnamese translation for OI Public Library
% Source-PDF: IOI中国国家候选队论文/国家集训队2022论文集.pdf
\documentclass[11pt]{article}
\usepackage{oipl-vn}

\title{Tập luận văn đội tuyển quốc gia Trung Quốc năm 2022}
\author{Tập luận văn đội tuyển quốc gia Olympic Tin học}
\date{China Computer Federation, tháng 2 năm 2022}

\begin{document}
\maketitle

\origtitle{Tập luận văn đội tuyển quốc gia Trung Quốc Olympic Tin học năm 2022}
\sourcepath{IOI China candidate team papers / National training team 2022 collection.pdf}

\renewcommand{\abstractname}{Tóm tắt}
\begin{abstract}
PDF nguồn là tập hợp luận văn đội tuyển quốc gia Trung Quốc năm 2022. Tệp này
là bản quét ảnh: \texttt{pdftotext} chỉ trả về các dấu sang trang và không khôi
phục được văn bản thân bài. Vì vậy bản LaTeX này chuyển ngữ phần nhận diện tập
sách, mục lục và hai bài đầu tiên bằng cách đọc trực tiếp các trang đã render
từ PDF, kết hợp OCR bằng \texttt{tesseract}.
\end{abstract}

\section{Thông tin đầu sách}

Tập sách có tiêu đề \emph{Tập luận văn đội tuyển quốc gia Trung Quốc Olympic
Tin học năm 2022}. Trang bìa ghi đơn vị China Computer Federation và thời điểm
phát hành là tháng 2 năm 2022. Tệp PDF gồm 246 trang; phần thân bắt đầu từ bài
về duy trì thông tin đường đi trên một lớp DAG đặc biệt.

\section{Mục lục đã dịch}

\begin{center}
\small
\begin{tabular}{r p{0.52\linewidth} p{0.24\linewidth}}
\textbf{Trang} & \textbf{Bài viết} & \textbf{Tác giả} \\
\hline
1 & Bàn về duy trì thông tin đường đi trên một lớp DAG đặc biệt & Dai Jiangqi \\
13 & DFT tổng quát hơn & Zhou Hangrui \\
28 & Bàn về ứng dụng của một tư tưởng tham lam dựa trên so sánh trong thi tin học & Zhang Junkai \\
49 & Bàn về các thuật toán liên quan tới đồ thị phẳng & Tang Shaoxuan \\
72 & Bàn về bài toán tương tác dạng tìm tham số tương tác & Hu Hao \\
82 & Bàn về thuật toán xấp xỉ trong OI & Xu Tingqiang \\
99 & Khảo sát bước đầu về bài toán tô màu đồ thị & Peng Bo \\
118 & Bàn về bài toán tối ưu không gian trong thi tin học & Chen Zhixuan \\
129 & Bàn về một lớp ma trận Monge & Zhang Jingxing \\
145 & Bàn về phân tích độ phức tạp của một lớp bảng băm & Chen Yusi \\
159 & Ứng dụng của đại số tuyến tính trong OI & Chang Ruinian \\
182 & Bàn về các bài toán tổ hợp chuỗi liên quan tới lý thuyết Lyndon & Wan Chengzhang \\
204 & Báo cáo ra đề \emph{Chuỗi của Cutie Duo} & Feng Shiyuan \\
213 & Bàn về bài toán mô phỏng luồng chi phí trong OI & Chen Kuowen
\end{tabular}
\end{center}

\section{Bàn về duy trì thông tin đường đi trên một lớp DAG đặc biệt}

\begin{center}
\textbf{Dai Jiangqi, Trường Ngoại ngữ Nam Kinh}
\end{center}

\subsection*{Tóm tắt}

Trong OI, không ít bài toán liên quan tới cấu trúc dữ liệu có thể được xem như
bài toán đường đi trên DAG. Bài viết này giới thiệu ba thuật toán có thể duy
trì thông tin đường đi trên DAG trong thời gian $O(n\operatorname{poly}(\log
S))$, trong đó $n$ là kích thước đầu vào, $S$ là tổng số đường đi của DAG và
$\operatorname{poly}(x)$ biểu thị một đa thức theo $x$. Bài viết cũng khảo sát
ứng dụng của các thuật toán này trong cấu trúc dữ liệu, chuỗi và số học.

\subsection*{Mở đầu}

Trong OI, rất nhiều bài toán có thể quy về mô hình liên quan tới đường đi trên
DAG, bao gồm nhưng không giới hạn ở đồ thị, cấu trúc dữ liệu và chuỗi. Trên DAG
tổng quát, việc duy trì nhanh thông tin đường đi là không thực tế, vì chỉ riêng
mô tả một đường đi đã cần $\Omega(n)$ bit. Do đó bài viết tập trung vào các DAG
có tổng số đường đi tương đối nhỏ; nhiều mô hình DAG thực dụng, chẳng hạn máy tự
động hậu tố, thỏa tính chất này.

Bài viết chủ yếu giới thiệu ba thuật toán có thể duy trì thông tin đường đi
trên DAG trong thời gian $O(\operatorname{polylog} S)$, với $S$ là tổng số đường
đi: phân rã chuỗi trên DAG, cây cân bằng bền vững và tách--trộn cây cân bằng.
Trong đó, phân rã chuỗi trên DAG có cài đặt gọn, thao tác sửa đổi tương đối dễ,
nên thường thực dụng hơn trong OI. Cây cân bằng bền vững có độ phức tạp thời
gian tối ưu và khả năng mở rộng mạnh, nhưng tốn bộ nhớ và khó cài đặt hơn. Độ
phức tạp thời gian của phương pháp tách--trộn cây cân bằng vẫn còn mở; hiện chỉ
có chứng minh cận trên $O(\log S\log q)$ mà chưa xây dựng được dữ liệu làm chặt
cận này. Tuy vậy, vì quan hệ giữa cập nhật và truy vấn của nó ngược với hai thuật
toán trước, tức truy vấn là offline còn DAG có thể bị bắt buộc đưa ra online,
nó vẫn có một số ứng dụng khó thay thế.

Phần 1 trình bày mô hình cơ bản của duy trì thông tin đường đi trên DAG. Phần 2
lần lượt trình bày ba thuật toán: phân rã chuỗi trên DAG, cây cân bằng bền vững
và tách--trộn cây cân bằng. Phần 3 nêu các ứng dụng trong một số bối cảnh.

\subsection{Mô hình cơ bản}

Cho một DAG có $n$ đỉnh, trong đó các cạnh ra của mỗi đỉnh có thứ tự. Sắp các
đỉnh mà đỉnh $i$ ($1\le i\le n$) trỏ tới thành dãy $G_i$. Đỉnh $i$ có trọng số
nguyên dương $c_i$, và cạnh ra thứ $j$ của nó có trọng số cạnh $j$.

Sắp tất cả các đường đi bắt đầu tại $1$ và kết thúc ở một đỉnh có bậc ra bằng
$0$ theo thứ tự từ điển của dãy trọng số cạnh; cách sắp này là duy nhất. Ký hiệu
dãy đường đi là $P_i$ ($1\le i\le S$), trong đó $S$ là tổng số đường đi thỏa điều
kiện.

Có $q$ truy vấn. Mỗi truy vấn cho một số nguyên dương $k$, yêu cầu tính tổng
trọng số đỉnh trên $P_k$.

Giả sử số cạnh của DAG không vượt quá $2n$. Dòng ràng buộc trong bản quét bị OCR
nhiễu; các phần đọc được gồm $k\le 10^{18}$ và các chặn đa thức thông thường cho
$n,q,c_i$.

\subsection{Lời giải}

Vì $k\le 10^{18}$, chỉ cần giữ lại $10^{18}$ đường đi đầu của DAG. Trên thực tế,
có thể dùng DP $O(n)$ để tính số đường đi từ mỗi đỉnh tới các đỉnh bậc ra $0$,
đồng thời lấy $\min$ với $10^{18}$, rồi chỉ giữ các đỉnh có giá trị DP nhỏ hơn
$10^{18}$. Ngoài ra cần sao chép các đỉnh nằm trên đường đi nhỏ thứ
$10^{18}$ theo thứ tự từ điển và giữ với các cạnh ra của chúng một tiền tố trỏ
tới nhóm đỉnh trước. Vì vậy chỉ cần xét trường hợp $S\le 10^{18}$.

Bằng cách nhị phân hóa bậc ra, dễ thấy ta chỉ cần xét trường hợp bậc ra của mỗi
đỉnh là $0$ hoặc $2$. Trong phân tích dưới đây giả sử $S>n$.

\subsubsection{Phân rã chuỗi trên DAG}

Nếu DAG là một cây có hướng, hoặc một rừng có hướng, tức bậc vào của mỗi đỉnh
không vượt quá $1$, hiển nhiên có thể dùng phân rã nặng--nhẹ để duy trì thông
tin trên chuỗi.

Ta mở rộng cách làm này lên DAG tổng quát hơn. Với mỗi cạnh $u\to v$, xét số
đường đi và số tiền tố đường đi liên quan tới hai đầu cạnh, rồi chọn cạnh nặng
theo đỉnh kế tiếp và tiền nhiệm có giá trị lớn nhất. Bản quét nhận diện nhiễu
ở công thức định nghĩa chính xác của $f_i,h_i,\operatorname{pre}_i$; phần đọc
được cho thấy $f_i$ là số đường đi kết thúc ở $i$, $h_i$ là số đường đi bắt đầu
từ $i$, còn $\operatorname{pre}_i$ là tiền nhiệm nặng của $i$.

Cạnh $u\to v$ là cạnh nặng khi $v$ là lựa chọn nặng của $u$ và $u$ là tiền nhiệm
nặng của $v$. Dễ thấy các cạnh nặng tạo thành một số chuỗi rời nhau, gọi là
chuỗi nặng.

Xét một đường đi bất kỳ trên DAG. Dọc đường đi đó, $f_i$ giảm dần còn $h_i$ tăng
dần. Với mỗi cạnh nhẹ, ít nhất một trong hai đại lượng phải giảm hoặc tăng theo
hệ số $2$. Vì vậy trên một đường đi chỉ có $O(\log S)$ cạnh nhẹ.

Do đó, từ đỉnh $1$ ta có thể mỗi lần tìm bằng nhị phân phần chung giữa đường đi
được hỏi và chuỗi nặng hiện tại; trên mỗi chuỗi nặng tiền xử lý tổng tiền tố.
Nhờ vậy bài toán ở trên được giải trong thời gian $O(n+q\log S\log n)$ và bộ nhớ
$O(n)$. Nút thắt nằm ở phép nhị phân trong mảng $f$.

\paragraph{Tối ưu.}
Tham khảo kỹ thuật cây nhị phân cân bằng toàn cục trong phân rã cây, độ phức tạp
có thể hạ xuống $O(n+q\log S)$. Với một chuỗi nặng cố định, đặt
$w_i=f_i-f_{i+1}$, với quy ước $f_{s+1}=0$. Dựng một cây nhị phân trên toàn bộ
chuỗi sao cho trong một đoạn $[l,r]$, gốc con là vị trí đầu tiên có tổng trọng số
hai phía đều không vượt quá một nửa tổng đoạn; nếu không tồn tại thì lấy đầu
đoạn. Tiền xử lý LCA trên cây nhị phân này.

Nếu $u$ là vị trí vào chuỗi nặng hiện tại và $v$ là vị trí ra, ta tiền xử lý LCA
giữa mỗi đỉnh với cuối chuỗi nặng rồi nhị phân từ $\operatorname{LCA}(u,k)$ đi
xuống. Độ phức tạp của phép nhị phân trên cây nhị phân chính là khoảng cách giữa
$\operatorname{LCA}(u,k)$ và $v$ trên cây đó.

Bắt đầu từ $v$ và đi lên hai bước mỗi lần, tổng trọng số của cây con ít nhất
tăng gấp đôi, trừ tình huống lá biên. Sau một bước đi xuống từ LCA, tổng cây con
không vượt quá đại lượng tương ứng ở điểm vào. Cộng trên mọi chuỗi nặng của một
truy vấn, vì đại lượng của chuỗi hiện tại không nhỏ hơn chuỗi kế tiếp, tổng độ
phức tạp nhị phân là $O(\log S)$.

Dùng splay hoặc treap cũng đạt cùng bậc độ phức tạp; bài viết không nhắc lại.

\paragraph{Mở rộng tối ưu.}
Dùng chia để trị, tức thông tin nửa nhóm trên đoạn với tiền xử lý
$O(n\log n)$ và truy vấn $O(1)$ dựa trên hợp nhất hai đoạn, kết hợp DP, có thể
hỗ trợ truy vấn giá trị lớn nhất của tổng trọng số đỉnh trên đoạn
$P_l,P_{l+1},\ldots,P_r$ trong thời gian $O(n\log n+q\log S)$. Cụ thể, với mỗi
đỉnh tính giá trị đường đi lớn nhất bắt đầu ở nó, rồi trên mỗi chuỗi nặng lưu cực
trị đoạn của các thông tin đó.

Với các bài toán có sửa điểm đơn, sửa chuỗi đơn giản và truy vấn thông tin trên
một chuỗi, cũng có thể tổng quát hóa cây đoạn có trọng số để giảm độ phức tạp
mỗi thao tác xuống $O(\log S)$.

Bản quét chứa một công thức chọn gốc cây nhị phân phụ thuộc vào hai dãy
$w_i=f_i-f_{i+1}$ và $x_i=h_i-h_{i-1}$, cùng tính chẵn lẻ của độ sâu $d$; công
thức này quá nhiễu để chép chính xác. Ý chính đọc được là khoảng cách từ điểm
vào tới LCA và từ điểm ra tới LCA đều bị chặn bởi các logarit của tỉ lệ tổng
trọng số, nên cộng trên mọi chuỗi nặng vẫn không vượt quá $O(\log S)$.

Ngoài ra, với thông tin khả nghịch có sửa điểm đơn, xây cây đoạn có trọng số
riêng theo hai phía cũng đạt truy vấn đơn $O(\log S)$. Tuy nhiên hằng số lớn; nếu
thông tin cần duy trì đơn giản thì thuật toán này không có ưu thế rõ rệt so với
BIT.

\subsubsection{Cây cân bằng bền vững}

Xét bài toán gốc theo thứ tự tô-pô đảo của DAG. Với mỗi đỉnh $i$, xét dãy $Q_i$
gồm mọi đường đi bắt đầu từ $i$.

Rõ ràng, nếu bậc ra của $i$ bằng $0$ thì $Q_i=((i))$. Nếu bậc ra của $i$ bằng
$2$, với hai cạnh ra theo thứ tự, $Q_i$ nhận được bằng cách ghép hai dãy đường
đi của hai con, rồi chèn $i$ vào đầu mỗi phần tử.

Ta chỉ duy trì tổng trọng số đỉnh của từng đường đi trong $Q_i$. Khi đó cần hỗ
trợ các thao tác sau trong $O(\log S)$:
\begin{enumerate}
  \item ghép hai dãy số nguyên độ dài không vượt quá $S$ thành một dãy mới;
  \item cộng cùng một hằng số vào mọi phần tử của một dãy;
  \item cho $k$, hỏi phần tử thứ $k$ của một dãy.
\end{enumerate}

Dùng cây cân bằng để duy trì. Do có thao tác tự sao chép, ta cần một loại cây
không có con trỏ cha và chiều cao chặt $O(\log S)$. Bài viết dùng WBLT, tức
Weight Balanced Leafy Tree. Trong WBLT, mỗi lá biểu diễn một phần tử; mỗi nút
trong có cây con trái và phải, và kích thước mỗi cây con không nhỏ hơn một hằng
số cân bằng nhân với tổng kích thước. Mỗi nút có thể lưu thông tin đoạn. Ở đây,
lá biểu diễn đường đi và mỗi nút giữ một nhãn cộng lười cho cây con.

Với thao tác ghép, chỉ cần merge hai WBLT. Khi ghép hai cây $A,B$ và giả sử
$|A|\ge |B|$, nếu tỉ lệ kích thước đã cân bằng thì tạo một gốc mới với hai cây
con $A,B$. Nếu không, xét cây con phải của $A$; khi ghép trực tiếp vẫn thỏa cân
bằng thì gắn $B$ vào phía phải. Nếu vẫn chưa thỏa, tiếp tục tách cấu trúc bên
phải của $A$, ghép các phần tương ứng rồi hợp nhất trở lại. Bằng quy nạp, ghép
hai WBLT có tỉ lệ kích thước là hằng số tốn $O(1)$; vì các đệ quy còn lại chỉ đi
về một phía, tổng thời gian ghép không vượt quá $O(\log S)$. Trong bài này còn
có thể lợi dụng tính khả trừ và giao hoán của thông tin để làm nhãn gần như
vĩnh viễn, giảm hằng số số lần \emph{pushdown}.

Thao tác cộng toàn dãy chỉ sửa nhãn ở gốc. Thao tác hỏi phần tử thứ $k$ đệ quy
từ gốc xuống lá. Vì vậy bài toán ban đầu được giải trong thời gian
$O(n+q\log S)$ và bộ nhớ $O(n\log S)$.

\paragraph{Mở rộng của lời giải WBLT.}
Ngoài ghép, WBLT và một số cây cân bằng khác còn hỗ trợ tách: cho cây $A$ và
$k\in[0,|A|]$, trả về hai WBLT lần lượt biểu diễn $k$ phần tử đầu và
$|A|-k$ phần tử sau.

Cụ thể, đệ quy từ trên xuống sẽ tách $A$ thành một số cây con có độ sâu tăng dần
và một số cây con có độ sâu giảm dần; sau đó ghép chúng từ dưới lên để thu được
hai cây kết quả. Vì chi phí ghép tỉ lệ với logarit tỉ lệ kích thước của hai
phần, dễ chứng minh thao tác tách tốn $O(\log S)$.

Nói cách khác, ngoài ghép dãy, sửa toàn cục và hỏi điểm đơn, cây cân bằng bền
vững còn hỗ trợ cắt lấy một đoạn dãy. Tất nhiên cũng có thể hỗ trợ sửa đoạn và
hỏi đoạn; trong bài toán gốc, đoạn này tương ứng với một đoạn của dãy đường đi
$P$.

Do cấu trúc phức tạp, cây cân bằng bền vững khó hỗ trợ sửa trọng số đỉnh của
DAG.

\paragraph{Thảo luận liên quan.}
Có thể thấy mỗi DAG mà mọi đỉnh có bậc ra $0$ hoặc $2$ tương đương về cấu trúc
với một cây Leafy Tree bền vững.

Điều này có nghĩa: khi không có tính chất đặc biệt và số thao tác ít, ta có thể
duy trì trực tiếp DAG để cài ghép dãy và cắt dãy, bỏ qua các thao tác giữ cân
bằng. Định nghĩa độ sâu $\operatorname{dep}_v$ của đỉnh $v$ là số cạnh trên
đường đi dài nhất bắt đầu ở $v$. Khi ghép $A,B$ thành $C$, chỉ cần tạo một đỉnh
mới có hai con $A,B$ và đặt
$\operatorname{dep}_C=\max(\operatorname{dep}_A,\operatorname{dep}_B)+1$.
Khi tách $A$ thành $B,C$, thời gian không vượt quá $\operatorname{dep}_A$, vì
trong quá trình đệ quy mỗi độ sâu chỉ bị truy cập ở trước và sau nhiều nhất một
lần. Bằng cách ghép từ sau ra trước dọc ranh giới giữa $B$ và $C$, ta đảm bảo
$\max(\operatorname{dep}_B,\operatorname{dep}_C)\le \operatorname{dep}_A$.
Sau $O(q)$ thao tác ghép và tách, độ sâu không vượt quá $q$, nên tổng độ phức tạp
là $O(q^2)$. Đây cũng đạt cận dưới lý thuyết trong mô hình này, vì sau $q$ bước
số đường đi có thể đạt $\exp(\Theta(q))$. Ưu thế của cây cân bằng bền vững nằm ở
các trường hợp DAG có cấu trúc đặc biệt hơn.

Ngoài ra, cấu trúc phân rã chuỗi trên DAG kết hợp cây đoạn có trọng số ở phần
trước cũng có thể được xem như một DAG ``cân bằng hơn'' nhưng tương đương với
DAG gốc theo ý nghĩa dãy đường đi.

\subsubsection{Tách--trộn cây cân bằng}

Xét phiên bản offline của bài toán gốc: tại đỉnh $1$, duy trì tập các chỉ số
$k$ xuất hiện trong mọi truy vấn.

Duyệt các đỉnh theo thứ tự tô-pô từ trước ra sau. Khi tới đỉnh $i$, mọi truy vấn
đang nằm tại $i$ được cộng thêm $c_i$ vào đáp án. Nếu $i$ có bậc ra $2$, chuyển
toàn bộ truy vấn này tới các đỉnh mà $i$ trỏ tới; nửa sau cần trừ đi kích thước
của nửa trước khỏi chỉ số $k$ tương ứng.

Cần duy trì nhiều dãy không giảm và hỗ trợ ba thao tác sau; sau khi bị thao tác,
dãy cũ biến mất:
\begin{enumerate}
  \item tách một dãy thành hai dãy mới theo ngưỡng $k$;
  \item cộng $k$ vào mọi phần tử của một dãy;
  \item trộn hai dãy đã sắp xếp thành một dãy mới.
\end{enumerate}

Dùng cây cân bằng để duy trì các dãy này thì hai thao tác đầu là trực tiếp.

Với thao tác trộn hai dãy $A$ và $B$ thành $C$, chia các vị trí
$1,2,\ldots,|A|+|B|$ thành $t$ đoạn liên tiếp cực đại sao cho các phần tử trong
mỗi đoạn của $C$ trước khi trộn đều thuộc cùng một dãy cũ. Dùng nhị phân từ trái
sang phải hoặc đệ quy từ gốc xuống đều cho lời giải đơn giản $O(t\log q)$.

Tiếp theo chứng minh tổng $t$ không vượt quá $O((n+q)\log S)$. Với một dãy $A$,
định nghĩa thế năng
\[
  \Phi(A)=\sum_i \log(2+A_{i+1}-A_i).
\]
Tổng thế năng luôn không âm và không vượt quá $O(q\log S)$. Thao tác tách làm
thế năng thay đổi nhiều nhất $O(\log S)$, thao tác cộng toàn dãy không đổi thế
năng. Vì vậy chỉ cần chứng minh trong một thao tác trộn, thế năng giảm ít nhất
$\Omega(t)-O(\log S)$.

Bản quét của chuỗi bất đẳng thức chi tiết bị nhiễu ở chỉ số, nhưng lập luận đọc
được như sau: gọi các đoạn cực đại sau trộn là $[l_i,r_i]$ và đặt
$d_i=l_i-r_{i-1}$ giữa hai đoạn kề nhau. Khi ghép hai đoạn kề, số hạng logarit
tương ứng được thay bằng logarit của khoảng cách lớn hơn; dùng
$\min(d_i,d_{i+1})$ và $\max(d_i,d_{i+1})$ để chặn, mỗi lần đổi đoạn đóng góp
một lượng giảm hằng số, còn hai biên chỉ tạo sai số $O(\log S)$. Do đó tổng số
đoạn cực đại trên mọi lần trộn bị chặn bởi tổng thế năng tăng.

Ta thu được lời giải thời gian $O((n+q)\log S\log q)$ và bộ nhớ $O(n+q)$.

\subsubsection{Mở rộng và thảo luận}

Nhìn lại bốn thao tác mà cây cân bằng bền vững ở phần trên hỗ trợ: ghép dãy, cắt
dãy, cộng toàn cục và hỏi điểm đơn. Nếu xem các thao tác này như một DAG, rồi
duyệt ngược và dùng cây cân bằng để duy trì mọi truy vấn, nút thắt vẫn là tách
cây, cộng toàn cục và trộn cây. Nói cách khác, vấn đề mà cây cân bằng bền vững
và vấn đề mà tách--trộn cây cân bằng giải quyết phần nào là chuyển vị của nhau.

Đáng tiếc là cận trên tốt nhất đã biết của tách--trộn cây cân bằng là
$O(\log S\log q)$, dù cũng chưa có dữ liệu đạt chặt cận này, còn cây cân bằng
bền vững đạt $O(\log S)$. Trong các bài toán liên quan tới bài viết này, thuật
toán tách--trộn chỉ hơn cây cân bằng bền vững khi có yêu cầu đặc biệt về bộ nhớ,
hoặc khi một số thông tin trong DAG thao tác bị bắt buộc đưa ra online, chẳng
hạn khi thao tác tách không phải tách theo ``$\ge k$'' và ``$<k$'' mà là tách
theo $k$ phần tử đầu và phần còn lại.

\subsection{Ứng dụng}

\subsubsection{Duy trì quá trình thao tác trên dãy}

Theo phần 2.2.2, xem DAG như một Leafy Tree bền vững giúp cài đặt thuận tiện
các thao tác cắt dãy và ghép dãy.

\paragraph{Ví dụ 1.}
Có $n$ dãy. Ban đầu dãy thứ $i$ chỉ chứa một số.

Cần thực hiện $q$ thao tác, mỗi thao tác thuộc một trong ba loại:
\begin{enumerate}
  \item tạo một dãy mới bằng cách ghép hai dãy đã có;
  \item tạo hai dãy mới, lần lượt là $k$ phần tử đầu của một dãy đã có và các
  phần tử còn lại, trong đó $k$ là số nguyên lớn không vượt quá độ dài dãy;
  \item xuất số nghịch thế của một dãy modulo $998244353$.
\end{enumerate}

Ràng buộc đọc được: $1<n<100$, $1<q<600$.

\paragraph{Lời giải ví dụ 1.}
Nếu không có thao tác loại 2, hiển nhiên có thể duy trì số lần xuất hiện của mỗi
giá trị trong từng dãy và tổng số nghịch thế, mỗi thao tác tốn $O(n)$.

Dùng cách làm ở phần 2.2.2: với thao tác loại 2, ta đệ quy từ trên xuống để tách
rồi ghép dần từ dưới lên, trong quá trình duy trì kích thước mỗi nút. Khi đó mỗi
thao tác tách có thể viết lại thành $O(q)$ thao tác ghép, còn mỗi thao tác ghép
tốn $O(n^2+n\omega)$, trong đó $\omega$ là số bit của kích thước, thường xem là
$32$ hoặc $64$. Tổng độ phức tạp là $O(n^2q^2+n\omega q^2)$.

\subsubsection{Ứng dụng trên máy tự động hậu tố}

Trong lý thuyết chuỗi của OI, máy tự động hậu tố (SAM) là một thuật toán hiệu
quả và thực dụng.

Cụ thể, với chuỗi $S$ trên bảng chữ cái $\Sigma$, xét tập vị trí kết thúc của
xâu con $T$ trong $S$, ký hiệu là $R_T$. Với các $T$ khác nhau, quan hệ giữa các
$R_T$ chỉ có thể là bao hàm hoặc rời nhau; vì vậy chỉ có $O(|S|)$ tập $R_T$ khác
nhau về bản chất và chúng tạo thành một cấu trúc cây, thực chất là cây hậu tố
của chuỗi đảo sau khi nén các đỉnh bậc hai.

Định nghĩa tự động $A$ có tập đỉnh là các $R_T$, với $T$ là xâu con của $S$, và
có cạnh ký tự $a$ từ $R_T$ tới $R_{Ta}$ nếu $Ta$ cũng là xâu con của $S$. Từ các
tính chất của SAM, $A$ có kích thước $O(|S|)$ và có thể xây dựng cây cùng tự động
trong $O(|S|)$.

Các xâu con khác nhau của $S$ tương ứng một-một với các đường đi bắt đầu từ xâu
rỗng trong $A$. Nhờ đó, thông tin về xâu con khác nhau được chuyển thành thông
tin đường đi trên tự động, giúp xử lý hiệu quả một số bài toán khá khó nếu chỉ
nhìn từ góc độ cây hậu tố.

\paragraph{Ví dụ 2.}
Nguồn bài là UOJ Long Round \#1, bài D\footnote{\url{https://uoj.ac/problem/577}},
đã lược bỏ chi tiết không cần thiết.

Cho chuỗi $S$ và mảng $c_i$. Với xâu con $T$ của $S$, ký hiệu tập vị trí kết
thúc của $T$ trong $S$ là $R_T$ và định nghĩa $f(T)=|R_T|\sum_{i\in R_T} c_i$.

Có $q$ truy vấn, mỗi truy vấn yêu cầu tính tổng $f(T)$ trên một họ xâu con liên
tiếp theo mô tả của đề gốc. Ràng buộc đọc được: $|S|\le 5\cdot 10^5$, bảng chữ
cái $26$. Một phần dòng truy vấn trong bản quét bị OCR nhiễu nên không chép lại
chính xác.

\paragraph{Lời giải ví dụ 2.}
Xây cây hậu tố của chuỗi đảo và SAM của $S$. Trên mỗi đỉnh của SAM, giá trị
$f$ là cố định, còn truy vấn đúng là tính tổng $f$ trên một đường đi của SAM.

Dùng phân rã chuỗi trên DAG và duy trì tổng tiền tố trên mỗi chuỗi nặng là giải
được. Nút thắt là tìm độ dài phần chung giữa đường đi truy vấn và chuỗi nặng
hiện tại; dùng bảng thưa LCP với tiền xử lý $O(|S|\log |S|)$ và truy vấn $O(1)$.

\paragraph{Thảo luận ví dụ 2.}
Thực ra SAM chỉ có $O(|S|)$ đường đi tối đại, tương ứng với mọi hậu tố của $S$.
Dựa trên tính chất này, tồn tại thuật toán đơn giản hơn.

Khi trải mọi đường đi của SAM ra, ta thu được một cây chỉ có $O(|S|)$ lá, thực
chất là cây hậu tố của chuỗi thuận. Mỗi truy vấn tương đương tính tổng trọng số
đỉnh trên một chuỗi của cây này; mặt khác các đỉnh này đều tương ứng với đỉnh
SAM.

Nén các đỉnh bậc hai của cây này và xét thao tác nén trên tự động: xóa mọi đỉnh
có bậc ra $1$ trên tự động, với mọi cạnh trỏ tới chúng thì đổi đích cạnh thành
đỉnh mà đích cũ trỏ tới, lặp đến khi đích có bậc ra khác $1$, đồng thời đổi ký
tự trên cạnh thành một chuỗi. Bằng cách lưu mảng trên cạnh, tự động sau nén vẫn
duy trì được tổng đoạn; khi trải tự động ra sẽ trực tiếp thu được một cây hậu tố
chuỗi thuận đã nén có $O(|S|)$ đỉnh. Xử lý tổng tiền tố trực tiếp trên cây này
là đủ. Tổng độ phức tạp là $O(|S|+q)$.

Đáng chú ý, tập đỉnh của SAM nén của chuỗi thuận và chuỗi đảo là giống nhau về
bản chất. Điều này đem lại một góc nhìn mới cho các bài toán xâu con liên quan
tới cả đầu trái và đầu phải. Tuy nhiên thuật toán này cũng có hạn chế: cấu trúc
SAM nén phụ thuộc vào tính đối xứng giữa chuỗi thuận và chuỗi đảo, nên khả năng
mở rộng kém hơn phân rã chuỗi trên DAG.

\paragraph{Ví dụ 3.}
Cho một trie $S$ có $n$ đỉnh. Định nghĩa xâu con của $S$ là chuỗi tương ứng với
một đường đi từ trên xuống.

Có $q$ truy vấn. Mỗi truy vấn cho $k$, yêu cầu tìm xâu con khác nhau nhỏ thứ
$k$ của $S$ hoặc báo không tồn tại. Để tránh đầu ra quá lớn, chỉ cần xuất tổng
mã ASCII của mọi ký tự trong xâu đó.

Ràng buộc đọc được: $n,q\le 10^5$, $k\le 10^{18}$, bảng chữ cái $26$.

\paragraph{Lời giải ví dụ 3.}
Cần xây SAM tổng quát. Ở đây $S$ được định nghĩa là trie tạo bởi các chuỗi trên
đường đi từ mọi đỉnh về gốc, còn tập \emph{right} của SAM tương ứng một-một với
đỉnh của trie.

Phân tích độ phức tạp của thuật toán xây SAM trên chuỗi không còn áp dụng trực
tiếp, và số cạnh của SAM lúc này có thể đạt $n|\Sigma|$. Theo các tài liệu được
trích dẫn, BFS trực tiếp trên trie rồi dùng thuật toán SAM cho chuỗi là tuyến
tính theo kích thước kết quả, nhưng tác giả không biết chứng minh đầy đủ. Cũng
có thể dùng suffix array tổng quát bằng nhân đôi để tính cây, rồi từ đó tính
SAM, tổng độ phức tạp $O(n\log n+n|\Sigma|)$; nếu dùng mở rộng của DC3 trên cây
có thể bỏ thừa số $\log$.

Sau khi xây được SAM, bài toán trở thành mô hình kinh điển của phần 1, và cả ba
thuật toán ở phần 2 đều giải được. Chẳng hạn dùng phân rã chuỗi trên DAG thì
tổng thời gian là $O(n\log n+n|\Sigma|+q\log^2 n)$, hoặc với tối ưu là
$O(n\log n+n|\Sigma|+q\log n)$.

\paragraph{Ví dụ 4.}
Nguồn bài là UOJ tuyển chọn đội, bài C\footnote{\url{https://uoj.ac/problem/697}}.

Cho chuỗi $S$ độ dài $n$ và $m$ xâu con của $S$ làm mẫu. Có $q$ truy vấn; mỗi
truy vấn yêu cầu tổng số lần xuất hiện của mọi xâu mẫu trong một xâu con của
$S$. Ràng buộc đọc được: $n\le 4\cdot 10^5$, $m\le 10^5$, $q\le 10^5$, bảng chữ
cái $26$.

\paragraph{Lời giải ví dụ 4.}
Xây SAM của $S$ và cây hậu tố của chuỗi đảo, rồi phân rã chuỗi trên DAG của SAM.

Với một xâu con $T$ của $S$, số xâu mẫu trong các hậu tố của $T$ được quyết định
bởi đỉnh tương ứng với $T$ trên SAM và các xâu mẫu tương ứng. Hơn nữa, chỉ cần
lấy số xâu mẫu trên đường đi từ đỉnh đó về gốc trong cây hậu tố đảo, rồi trừ đi
số xâu mẫu có độ dài lớn hơn $|T|$.

Với mỗi truy vấn, định vị xâu con thành một đường đi trên tự động và tính tổng
đóng góp của mọi đỉnh trên đường đi đó. Định vị đường đi có thể làm trong
$O((n+q)\log n)$ bằng cách tiền xử lý $O(n\log n)$ trên chuỗi gốc và các chuỗi
nặng của DAG rồi hỏi LCP xâu con trong $O(1)$.

Với thống kê đáp án: loại đóng góp thứ nhất chỉ cần lưu tổng tiền tố của
$w_v$ trên mỗi chuỗi nặng; loại thứ hai chuyển thành bài toán đếm điểm trong
hình bình hành nghiêng $45^\circ$ trên từng chuỗi nặng, sau đổi tọa độ là bài
toán đếm điểm hai chiều thông thường. Tùy cách cài LCP và đếm điểm, tổng độ phức
tạp là $O((n+m)\log n+q\log^2 n)$ hoặc $O((n+m+q)\log n)$.

\subsubsection{Thuật toán Euclid vạn năng}

Trên lưới có đường thẳng $y=\frac abx$ với $a,b\in\mathbb{N}^{+}$. Xét một điểm
động trên đường thẳng này di chuyển theo hướng $x+y$ tăng trên một đoạn. Mỗi khi
đi qua một đường ngang, ghi ký hiệu $A$; mỗi khi đi qua một đường dọc, ghi ký
hiệu $B$; nếu đi qua điểm lưới thì ghi $A$ trước rồi $B$ sau. Cuối cùng thu được
một dãy trên $\{A,B\}$, ví dụ với $a=3,b=2$ thì dãy là $ABBAB$.

Thay $A$ và $B$ bằng hai phần tử của một nửa nhóm, tức thông tin ghép được trong
$O(1)$, mục tiêu của thuật toán Euclid vạn năng là tính phần tử nửa nhóm ứng với
toàn bộ dãy.

Trước hết giả sử đoạn di chuyển là từ $(0,0)$ tới $(a,b)$ và $\gcd(a,b)=1$. Nếu
$a>b$, đặt $k=\lfloor a/b\rfloor$. Dễ thấy sau mỗi $A$ đều có $k$ ký hiệu $B$.
Thay mọi đoạn $AB^k$ trong dãy gốc bằng một ký hiệu mới tương ứng; về mặt hình
học, điều này tương đương tác động một phép biến đổi tọa độ lên mặt phẳng và
đường thẳng, đưa bài toán về kích thước nhỏ hơn. Lặp quá trình $O(\log a)$ lần
sẽ giải được bài toán.

Dùng DAG để mô tả thuật toán trên: ban đầu có hai đỉnh $A,B$. Mỗi lần tạo một
đỉnh $C$, nối cạnh tới $A$ và $B$ theo phép thay thế tương ứng, rồi thay
$(a,b)$ bằng $(a-b,b)$ hoặc $(a,b-a)$ tùy bên lớn hơn, cho tới khi $a=b$.

Số cạnh của DAG có thể đạt độ phức tạp của phép trừ lặp, tức $O(a)$. Chỉ cần
dùng nhân đôi để tối ưu các thao tác trừ liên tiếp là có thể giảm kích thước đồ
thị xuống $O(\log(a+b))$.

Dãy $AB$ do quá trình Euclid vạn năng sinh ra chính là dãy các đỉnh cuối của
mọi đường đi khi trải DAG theo thứ tự DFS. Vì vậy phiên bản tổng quát của thuật
toán Euclid vạn năng tương đương truy vấn đoạn trên Leafy Tree bền vững ứng với
DAG này. Nhiều bài toán Euclid cổ điển, như tổng $\lfloor ai/b\rfloor$ hoặc
$\min_i(ai\bmod b)$, có thể chuyển thành mô hình này; hơn nữa DAG này cho phép
giải một số bài toán phức tạp hơn liên quan tới quá trình Euclid.

\paragraph{Ví dụ 5.}
Đây là bài mẫu thuật toán Euclid vạn năng trên LOJ\footnote{\url{https://loj.ac/p/6440}}.

Cho $n,m,l$ và hai ma trận vuông $A,B$ kích thước không vượt quá $10$. Cần tính
biểu thức ma trận theo dãy Euclid do đề bài định nghĩa. Bản quét của công thức
mục tiêu bị nhiễu nên không chép lại an toàn. Ràng buộc đọc được:
$n,m,l\le 10^{18}$.

\paragraph{Lời giải ví dụ 5.}
Xét dãy $AB$ ứng với đường thẳng liên quan trong đề. Duy trì ma trận $M$, ban đầu
$M=I$. Từ vị trí thứ hai của dãy trở đi, mỗi khi gặp $A$ thì đặt $M=MB$; mỗi khi
gặp $B$ thì cộng $M$ vào đáp án rồi đặt $M=AM$. Khi kết thúc, đáp án là giá trị
cần tìm.

Tác động của một đoạn của dãy có thể mô tả bằng một bộ ma trận: nếu trước khi đi
vào đoạn có trạng thái $M$, sau đoạn đáp án được cộng thêm một biểu thức tuyến
tính theo $M$ và trạng thái mới cũng là một biến đổi tuyến tính của $M$. Vì vậy
thông tin đoạn ghép được trong $O(1)$, và có thể áp dụng trực tiếp phiên bản
truy vấn đoạn của thuật toán Euclid vạn năng.

\paragraph{Ví dụ 6.}
Nguồn bài là CTT 2021 Day 4 T3, đã lược bỏ chi tiết không cần thiết.

Xét đồ thị vô hướng $n$ đỉnh, trong đó $i$ nối với $(i+d)\bmod n$. Trọng số đỉnh
lấy theo một dãy tuần hoàn độ dài $m$: $w_i=a_{i\bmod m}$.

Cần thực hiện $q$ lần sửa điểm đơn trên trọng số, sau mỗi lần sửa tính kích
thước tập độc lập có trọng số lớn nhất của đồ thị.

Ràng buộc đọc được: $d<n\le 10^9$, $m\le 2\cdot 10^5$, $w_i\le 400$,
$q\le 10^5$.

\paragraph{Lời giải ví dụ 6.}
Dễ thấy đồ thị gồm $\gcd(n,d)$ chu trình, và chỉ có một số loại chu trình khác
nhau về bản chất. Sau khi xóa mọi đỉnh từng bị sửa ít nhất một lần, đồ thị còn
lại gồm $O(q)$ chuỗi và một số loại chu trình. Chỉ cần tiền xử lý đáp án của
mỗi chuỗi và mỗi loại chu trình rồi cộng lại; với mỗi chuỗi, dùng một ma trận
$2\times2$ biểu diễn trạng thái chọn hoặc không chọn hai đầu để tính tập độc lập
có trọng số lớn nhất, rồi duy trì cây đoạn trên các điểm quan trọng.

Xét thông tin trên một chu trình. Trên lưới, vẽ đoạn $y=\frac dnx$ với
$x\in[0,n)$ và xây dãy cùng DAG tương ứng của thuật toán Euclid vạn năng. Đồng
thời duy trì một giá trị $x$; các giá trị $x$ khác nhau tương ứng với các chu
trình khác nhau. Duyệt dãy và duy trì ma trận $2\times2$; mỗi khi gặp một loại
ký hiệu thì cập nhật $x$ theo modulo $m$, đồng thời ghép thông tin tương ứng với
$w_x$. Sau khi thao tác xong, ma trận nhận được là đáp án của chu trình; nếu chỉ
thao tác trên một đoạn thì nhận được đáp án của chuỗi.

Thông tin đoạn của dãy có thể định nghĩa là $m$ ma trận $2\times2$, biểu diễn
tác động khi giá trị ban đầu của $x$ lần lượt là $0,1,\ldots,m-1$, cùng một độ
dời cho $x$ sau khi đi qua đoạn. DP trên DAG Euclid tính được mọi thông tin đoạn
trong $O(m\log n)$, từ đó lập tức có đáp án chu trình.

Đáp án chuỗi là truy vấn đoạn trên thông tin đường đi của DAG. Vì DAG này chỉ có
$O(\log n)$ đỉnh, DFS trực tiếp từ trên xuống đạt $O(q\log n)$. Tổng độ phức tạp
là $O((m+q)\log n)$.

\subsection*{Kết luận}

Bài viết đã giới thiệu ba thuật toán: phân rã chuỗi trên DAG, cây cân bằng bền
vững và tách--trộn cây cân bằng, đồng thời đưa ra ứng dụng của chúng trong một
số bài OI. Tác giả phỏng đoán phía sau mô hình duy trì đường đi trên DAG có thể
còn những tính chất đẹp và sâu hơn; các mở rộng khác của ba thuật toán và mối
liên hệ giữa chúng vẫn cần được khai thác thêm. Hy vọng bài viết có thể gợi mở
hướng nghiên cứu tiếp theo cho độc giả quan tâm.

\subsection*{Lời cảm ơn}

Tác giả cảm ơn China Computer Federation đã cung cấp nền tảng học tập và trao
đổi; cảm ơn huấn luyện viên Yang Yuliang của đội tuyển quốc gia đã hướng dẫn;
cảm ơn cha mẹ, nhà trường, các thầy cô và bạn học đã bồi dưỡng, giảng dạy và hỗ
trợ; cảm ơn các anh chị và bạn bè đã trao đổi, góp ý, truyền cảm hứng và đọc
bản thảo; cảm ơn mọi nguồn tư liệu tham khảo; và cảm ơn độc giả đã dành thời
gian đọc bài viết.

\subsection*{Tài liệu tham khảo}

\begin{enumerate}
  \item Zhang Zheyu. \emph{Bàn về thuật toán chia để trị trên cây}. Luận văn đội
  tuyển quốc gia IOI 2019, 2019.
  \item Wang Siqi. \emph{Leafy Tree và cây cân bằng trọng số được cài đặt từ nó}.
  Luận văn đội tuyển quốc gia IOI 2018, 2018.
  \item SF. \emph{Bàn về máy tự động hậu tố nén}. Luận văn đội tuyển quốc gia
  IOI 2020, 2020.
  \item Liu Yanyi. \emph{Mở rộng máy tự động hậu tố trên trie}. Luận văn đội
  tuyển quốc gia IOI 2015, 2015.
  \item J. A. Blumer. \emph{Algorithms for the directed acyclic word graph and
  related structures}. Ph.D. Thesis, Denver University, 1985.
  \item A. Blumer, J. Blumer, D. Haussler, R. M. McConnell, A. Ehrenfeucht.
  \emph{Complete inverted files for efficient text retrieval and analysis}.
  J. ACM 34(3)(1987), pp. 578--595.
  \item J. Karkkainen, P. Sanders, S. Burkhardt. \emph{Linear work suffix array
  construction}. J. ACM 53(6)(2006), pp. 918--936.
\end{enumerate}

\section{DFT tổng quát hơn}

\begin{center}
\textbf{Zhou Hangrui, Trường Văn Uyên thuộc Tập đoàn Giáo dục Trung học Học
Quân Hàng Châu}
\end{center}

\subsection*{Tóm tắt}

Biến đổi Fourier rời rạc được dùng rất rộng rãi trong thi tin học. Bài viết này
mở rộng DFT truyền thống, dùng độ phức tạp hơi cao hơn để thực hiện phép chập
do một quan hệ hai ngôi tổng quát hơn định nghĩa.

\subsection{Định nghĩa liên quan và quy ước}

\subsubsection{Phép chập và các phép toán cơ bản}

\paragraph{Định nghĩa 1.1.}
Với một tập $G$ và một quan hệ hai ngôi $\varphi:G\times G\to G$, cho hai ánh xạ
$A,B:G\to \mathbb C$. Định nghĩa phép chập $C=A\times B$ bởi
\[
  C(v)=\sum_{\varphi(x,y)=v} A(x)B(y).
\]

Tổng của hai ánh xạ $C=A+B$ thỏa
\[
  C(v)=A(v)+B(v).
\]

Tích của một ánh xạ với một số phức $B=cA$ thỏa
\[
  B(v)=cA(v).
\]

Bài viết chỉ xét trường hợp $(G,\varphi)$ tạo thành một nửa nhóm hữu hạn. Với
ánh xạ $A:G\to\mathbb C$, trong bài này $A_x$ và $A(x)$ có cùng ý nghĩa.

\subsubsection{DFT}

\paragraph{Định nghĩa 1.2.}
Với dãy kích thước $(a_i)_{i=1}^{k}$, định nghĩa tập dãy ma trận tương ứng là
$R$, thỏa:
\begin{itemize}
  \item $R=\{(m_i)_{i=1}^{k}\mid m_i\in M_{a_i}\}$, trong đó $M_n$ chứa mọi ma
  trận phức $n\times n$;
  \item $C=A+B$ thỏa $C_i=A_i+B_i$;
  \item $C=AB$ thỏa $C_i=A_iB_i$;
  \item $B=cA$ thỏa $B_i=cA_i$.
\end{itemize}
Ở đây $A,B,C\in R$ là các dãy ma trận có cùng dãy kích thước, còn $c$ là một số
phức.

\paragraph{Định nghĩa 1.3.}
Với một quan hệ chập, định nghĩa DFT tương ứng của nó là một biến đổi tuyến tính
khả nghịch
\[
  \mathcal F:(G\to\mathbb C)\to R,
\]
trong đó $R$ là tập dãy ma trận do một dãy kích thước $(a_i)_{i=1}^{k}$ định
nghĩa, đồng thời thỏa:
\begin{itemize}
  \item $\mathcal F(A)\times\mathcal F(B)=\mathcal F(AB)$;
  \item $\mathcal F(A)+\mathcal F(B)=\mathcal F(A+B)$;
  \item $c\mathcal F(A)=\mathcal F(cA)$;
  \item $\sum_{i=1}^{k}a_i^2=|G|$.
\end{itemize}
Có thể xem $\mathcal F$ là một đẳng cấu.

Với phép nhân dãy ma trận, ta làm được trong
$O(\sum_{i=1}^{k}a_i^\omega)$, trong đó $\omega$ là số sao cho phép nhân ma
trận $n\times n$ có độ phức tạp $O(n^\omega)$; hiện đã biết
$\omega\le 2.3728596$~[5].

Không gian $(G\to\mathbb C)$ có một hệ cơ sở $(A_x)(x\in G)$, trong đó
$A_x$ là ánh xạ
\[
  t\mapsto [t=x]=
  \begin{cases}
    1,& t=x,\\
    0,& t\ne x.
  \end{cases}
\]
Vì mọi phần tử đều phân rã được thành tổ hợp tuyến tính của cơ sở, ta chỉ cần
quan tâm tới $\mathcal F(A_x)$. Đặt $\mathcal F_x=\mathcal F(A_x)$.

\paragraph{Bổ đề 1.1.}
Một quan hệ DFT $\mathcal F$ tương ứng một-một với họ $(\mathcal F_x)$ thỏa:
\begin{itemize}
  \item $\mathcal F_x\times\mathcal F_y=\mathcal F_{\varphi(x,y)}$;
  \item các $\mathcal F_x$ tạo thành một hệ cơ sở của $R$.
\end{itemize}

\paragraph{Chứng minh.}
Với một quan hệ DFT $\mathcal F$, do tính đẳng cấu, họ $(\mathcal F_x)$ mà nó
thu được hiển nhiên hợp lệ.

Ngược lại, với một họ $(\mathcal F_x)$, có thể dựng
\[
  \mathcal F(A)=\sum_{x\in G} A(x)\mathcal F_x.
\]
Đối với phép nhân:
\[
\begin{aligned}
\mathcal F(A)\times\mathcal F(B)
&=\sum_{x\in G}\sum_{y\in G} A_xB_y\mathcal F_x\times\mathcal F_y\\
&=\sum_{v\in G}\sum_{\varphi(x,y)=v} A_xB_y\mathcal F_x\times\mathcal F_y\\
&=\sum_{v\in G}\mathcal F_v\sum_{\varphi(x,y)=v} A_xB_y\\
&=\mathcal F(A\times B).
\end{aligned}
\]
Phép cộng và nhân vô hướng hiển nhiên thỏa. Tính song ánh suy ra từ việc
$(\mathcal F_x)$ tạo thành một hệ cơ sở. \hfill\rule{0.55em}{0.55em}

Do đó, khi đã biết $(\mathcal F_x)$, DFT có thể xem như phép nhân một véc-tơ
hàng với ma trận biến đổi, độ phức tạp $O(|G|^2)$. Nếu cần IDFT, tức biến đổi
ngược của DFT, việc tìm ma trận nghịch đảo cần tiền xử lý $O(|G|^\omega)$.

\subsection{DFT truyền thống và nhóm cyclic}

\subsubsection{Phép chập dãy}

Với phép chập dãy thông thường, quan hệ hai ngôi tương ứng là $(\mathbb Z,+)$.
Tuy nhiên khi thao tác thực tế, ta thường lấy $2^k$ sao cho $2^k$ vượt quá độ
dài dãy $n$, rồi chuyển quan hệ hai ngôi thành $(\mathbb Z_{2^k},+)$; ở đây
$\mathbb Z_p$ biểu thị nhóm cyclic cấp $p$.

\subsubsection{Nhóm cyclic}

\paragraph{Bổ đề 2.1.}
Với nhóm cyclic $\mathbb Z_p$, một DFT tương ứng là
\[
  \mathcal F_i=\{[1],[\omega_p^i],[\omega_p^{2i}],\ldots,
  [\omega_p^{(p-1)i}]\},
\]
trong đó $\omega_p$ là một căn đơn vị nguyên thủy bậc $p$.

Độc giả có thể tự kiểm tra rằng cách dựng này thỏa các điều kiện trên.

\subsubsection{FFT và biến đổi chirp-Z}

Với nhóm cyclic $\mathbb Z_p$ thỏa $p=2^k$, có thể dùng FFT để tối ưu, đưa DFT
về $O(p\log p)$.

Với nhóm cyclic $\mathbb Z_p$ tổng quát hơn, có thể dùng biến đổi chirp-Z để
chuyển về trường hợp trên, cũng đạt $O(p\log p)$. Cụ thể:
\[
\begin{aligned}
\mathcal F(A)_i
&=\left[\sum_{j=0}^{p-1} A_j\omega_p^{ij}\right]\\
&=\left[\frac{1}{\omega_p^{\binom{i}{2}}}
  \sum_{j=0}^{p-1}\frac{A_j}{\omega_p^{\binom{j}{2}}}
  \omega_p^{\binom{i+j}{2}}\right].
\end{aligned}
\]
Đây là dạng của một phép chập trừ; sau khi đảo thứ tự có thể chuyển thành phép
chập dãy, rồi chuyển tiếp thành bài toán trên $\mathbb Z_{2^k}$. Độ phức tạp là
$O(p\log p)$.

Vì vậy với nhóm cyclic, ta đều có thể hoàn thành DFT trong
$O(|G|\log |G|)$; độ phức tạp của phép nhân ma trận hiển nhiên là $O(|G|)$.

\subsection{Tổng trực tiếp và DFT tương ứng}

\subsubsection{Tổng trực tiếp}

\paragraph{Định nghĩa 3.1.}
Với hai quan hệ hai ngôi $(G_1,\times)$ và $(G_2,\times)$, định nghĩa tổng trực
tiếp $(G,\times)$ của chúng bởi:
\begin{itemize}
  \item $G=\{(x,y)\mid x\in G_1,\ y\in G_2\}$;
  \item $(x_1,y_1)\times(x_2,y_2)=(x_1\times x_2,y_1\times y_2)$.
\end{itemize}

\subsubsection{Xây dựng DFT}

Trước hết chứng minh một bổ đề.

\paragraph{Bổ đề 3.1.}
Với không gian tuyến tính hữu hạn chiều, ký hiệu tích tensor là $\otimes$. Đặt
$V=V_1\otimes V_2=\operatorname{span}\{x\mid x=a\otimes b,\ a\in V_1,\ b\in
V_2\}$. Nếu một cơ sở của $V_1$ là $(a_i)_{i=1}^{n}$, một cơ sở của $V_2$ là
$(b_i)_{i=1}^{m}$, thì một cơ sở của $V$ là
\[
  c=\{x\mid x=a_i\otimes b_j,\ i\in[1,n],\ j\in[1,m]\}.
\]

\paragraph{Chứng minh.}
Hiển nhiên $c$ là một hệ sinh của $V$, nên chỉ cần chứng minh các phần tử của
nó độc lập tuyến tính. Phản chứng, đặt
$c_{(i-1)m+j}=a_i\otimes b_j$ và giả sử
\[
  \sum_{i=1}^{n}\sum_{j=1}^{m} c_{(i-1)m+j}k_{i,j}=0
\]
với không phải mọi $k_{i,j}$ đều bằng $0$. Đặt
$S_i=\sum_{j=1}^{m} k_{i,j}b_j$, ta có $\sum_{i=1}^{n} S_i\otimes a_i=0$.
Vì các $a_i$ tạo thành một cơ sở, suy ra $\forall i,\ S_i=0$. Vì các $b_j$ tạo
thành một cơ sở, nếu $S_i=0$ thì $k_{i,j}=0$ với mọi $j\in[1,m]$. Do đó mọi
$k_{i,j}$ đều bằng $0$, mâu thuẫn. \hfill\rule{0.55em}{0.55em}

\paragraph{Định nghĩa 3.2.}
Với hai dãy ma trận $R$ và $R'$ có dãy kích thước lần lượt là $(a_i)_{i=1}^{n}$
và $(a'_i)_{i=1}^{m}$, định nghĩa tích tensor của chúng $R\otimes R'$ bởi
\[
  (R\otimes R')_{(i-1)m+j}=R_i\otimes R'_j.
\]

\paragraph{Định lý 3.1.}
Nếu $G_1\to\mathbb C$ và $G_2\to\mathbb C$ đều có DFT, thì
$G\to\mathbb C$ cũng có DFT.

\paragraph{Chứng minh.}
Giả sử hai DFT lần lượt là $\mathcal A,\mathcal B$, với dãy kích thước tương ứng
$(a_i)_{i=1}^{n}$ và $(b_i)_{i=1}^{m}$. Giả sử DFT của $G\to\mathbb C$ là
$\mathcal F$, dãy kích thước tương ứng là $(c_i)_{i=1}^{nm}$. Có thể dựng như
sau:
\begin{itemize}
  \item $c_{(i-1)m+j}=a_ib_j$, với $i\in[1,n]$, $j\in[1,m]$;
  \item $\mathcal F_{(x,y),(i-1)m+j}=\mathcal A_{x,i}\otimes\mathcal B_{y,j}$,
  trong đó $\otimes$ là tích tensor.
\end{itemize}
Với phép cộng và nhân vô hướng, kết luận hiển nhiên.

Do tính chất của tích tensor
$(x_1\otimes y_1)(x_2\otimes y_2)=(x_1x_2)\otimes(y_1y_2)$, trong đó
$x_1,x_2$ là ma trận $n\times n$ và $y_1,y_2$ là ma trận $m\times m$, với phép
nhân ta có:
\[
\begin{aligned}
\mathcal F_{(x_1,y_1),v}\times\mathcal F_{(x_2,y_2),v}
&=(\mathcal A_{x_1,a}\otimes\mathcal B_{y_1,b})
  (\mathcal A_{x_2,a}\otimes\mathcal B_{y_2,b})\\
&=(\mathcal A_{x_1,a}\mathcal A_{x_2,a})\otimes
  (\mathcal B_{y_1,b}\mathcal B_{y_2,b})\\
&=\mathcal F_{(x_1,y_1)\times(x_2,y_2),v},
\end{aligned}
\]
trong đó $v=(a-1)m+b$, $a\in[1,n]$, $b\in[1,m]$.

Theo Bổ đề 3.1, các $(\mathcal F_x)$ tạo thành một cơ sở. \hfill\rule{0.55em}{0.55em}

Suy ra hệ quả sau.

\paragraph{Bổ đề 3.2.}
Nếu
\[
  G\to\mathbb C=(G_1\times G_2\times\cdots\times G_k)\to\mathbb C
\]
và mọi $G_i\to\mathbb C$ ($i\in[1,k]$) đều có DFT, thì có thể dựng DFT của
$G\to\mathbb C$.

\subsubsection{Độ phức tạp tốt hơn}

\paragraph{Bổ đề 3.3.}
Ma trận biến đổi ứng với $G$ là tích tensor của ma trận biến đổi ứng với $G_1$
và ma trận biến đổi ứng với $G_2$.

Điều này suy ra ngay từ quá trình dựng DFT của $G$.

\paragraph{Định lý 3.2.}
Với $G\to\mathbb C=(G_1\times G_2)\to\mathbb C$, ánh xạ
$a:G\to\mathbb C$, việc áp dụng DFT ứng với $G$ lên $a$ tương đương lần lượt áp
dụng DFT ứng với $G_1$ và $G_2$. Cụ thể:
\begin{itemize}
  \item Với $x\in G_1$, đặt $G'_x$ thỏa $G'_{x,y}=G_{(x,y)}$. Áp dụng DFT trên
  $G_2$ cho từng $G'_x$, ghi kết quả là $H_x$.
  \item Gọi tập dãy ma trận ứng với $G_2$ là $R$, dãy kích thước tương ứng là
  $(b_i)_{i=1}^{k}$. Với $i\in[1,k]$, đặt $H'_{i,x}=H_{x,i}$. Áp dụng DFT trên
  $G_1$ cho từng $H'_i$. Chú ý ánh xạ lúc này trở thành $G_1\to M_{b_i}$, nhưng
  quá trình tương tự.
\end{itemize}
Chỉ cần xét tính chất của ma trận biến đổi là đủ.

Tương tự, có hệ quả:

\paragraph{Bổ đề 3.4.}
Nếu
\[
  G\to\mathbb C=(G_1\times G_2\times\cdots\times G_k)\to\mathbb C,
\]
thì áp dụng DFT ứng với $G$ lên $a$ tương đương lần lượt áp dụng DFT ứng với
$G_i$ ($i\in[1,k]$), có thể cần chia lại $G$.

\paragraph{Bổ đề 3.5.}
Độ phức tạp của cách làm này là
\[
  T=O\!\left(|G|\sum_{i=1}^{k}\frac{T_i}{|G_i|}\right),
\]
trong đó $T_i$ biểu thị độ phức tạp DFT trên $G_i$.

Ta có các hệ quả:
\begin{itemize}
  \item nếu $T_i=O(|G_i|\log |G_i|)$, thì $T=O(|G|\log |G|)$;
  \item nếu $T_i=O(|G_i|^2)$, thì $T=O(|G|\sum |G_i|)$;
  \item nếu mọi $|G_i|$ đều là hằng số, có thể xem
  $T=O(k|G|)=O(|G|\log |G|)$.
\end{itemize}

Ví dụ, FWT $n$ chiều là tổng trực tiếp của $n$ bản sao $\mathbb Z_2\to\mathbb C$,
độ phức tạp là $O(|G|\log |G|)=O(n2^n)$.

\paragraph{Ví dụ 1 (Numbers).}
Nguồn bài: Zhang Huaihuai, \url{https://ioihw21.duck-ac.cn/problem/101}.

Cho $n,k$ và hai mảng độ dài $k$ là $m,p$ với $\sum_{i=1}^{k}p_i=1$. Trạng thái
là một $k$-bộ $c$, ban đầu mọi phần tử đều bằng $0$. Mỗi lần thao tác, với xác
suất $p_i$ ta đặt $c_i:=(c_i+1)\bmod m_i$. Hãy tính kỳ vọng của số trạng thái
khác nhau về bản chất đã đi qua trước lần đầu quay lại trạng thái ban đầu.
$T=\prod m_i\le 5\cdot 10^5$. Dữ liệu được đảm bảo ngẫu nhiên, không cần xét
các trường hợp biên.

\paragraph{Lời giải.}
Đặt $G=\mathbb Z_{m_1}\times\mathbb Z_{m_2}\times\cdots\times\mathbb Z_{m_k}$,
và $0$ là trạng thái ban đầu.

Xét việc tính xác suất mỗi trạng thái được đi qua; dưới đây giả sử cần tính xác
suất trạng thái $x\in G$ được đi qua. Giả định rằng khi đạt tới trạng thái ban
đầu hoặc $x$, thao tác lập tức kết thúc. Gọi $f_i$ là số lần kỳ vọng đi qua
trạng thái $i$ trước khi đạt tới trạng thái ban đầu hoặc $x$. Đặc biệt, cưỡng
bức đặt $f_0=1$, $f_x=0$. Với trạng thái khác $y$, ta có
\[
  f_y=\sum_{i=1}^{k} f_{d(y,i)}p_i,
\]
trong đó $d(y,i)$ biểu thị trạng thái thu được khi giảm phần tử thứ $i$ của
trạng thái $y$ đi $1$. Khử Gauss trực tiếp làm được trong $O(T^4)$.

Đưa vào biến $d_0,d_x$ thỏa:
\[
  f_0=\sum_{i=1}^{k} f_{d(0,i)}p_i-d_0=1,\qquad
  f_x=\sum_{i=1}^{k} f_{d(x,i)}p_i-d_x=0.
\]
Kiểm tra cho thấy $d_x$ chính là giá trị cần tìm.

Gọi $F$ là ánh xạ $x\mapsto f_x$, và $P$ là ánh xạ
\[
  x\mapsto
  \begin{cases}
    p_i,& x_j=1\Longleftrightarrow i=j,\\
    0,& \text{ngược lại}.
  \end{cases}
\]
Tương tự, với $d_0,d_x$ định nghĩa $D:y\mapsto d_y$. Khi đó các phương trình
trên có thể viết thành:
\[
  F=F\times P-D,\qquad f_0=1,\qquad f_x=0.
\]
Mặt khác, $F=F\times P-D$ tương đương
\[
  \mathcal F(F)=\mathcal F(F)\times\mathcal F(P)-\mathcal F(D),
\]
tức tương đương $T$ phương trình:
\[
  \mathcal F_i(P_i-1)=\mathcal D_i\qquad (i\in G).
\]
Ở đây $\mathcal F_g$ ($g_i\in G$) và $\mathcal F(F)_i$ ($i\in[1,|G|]$) tương
ứng một-một; $\mathcal P,\mathcal D$ được định nghĩa tương tự.

Ta thấy $\mathcal P_0=\sum p_i=1$; thực tế, với nhóm hữu hạn, mọi kết quả DFT
đều có cộng dồn. Phương trình tương ứng là
$\mathcal D_0=d_0+d_x=0$. Do dữ liệu ngẫu nhiên, xác suất để $\mathcal P_i$
($i\ne0$) bằng $0$ là rất nhỏ.

Định nghĩa $t(a,b)$ bởi
\[
  \mathcal F(C)_b=\sum C_a t(a,b),
\]
tức là hệ số của ma trận biến đổi. Khi đó
$t(a,b)=t(b,a)$ và $t(a,b)^{-1}=t(a^{-1},b)$. Ngoài ra, $f_0=1$ tương đương
$\sum\mathcal F_i=|G|$, còn $f_x=0$ tương đương
$\sum\mathcal F_i/t(i,x)=0$.

Lúc này hệ phương trình phức tạp ban đầu được chuyển thành dạng:
\[
  \sum \mathcal F_i=|G|, \tag{1}
\]
\[
  \sum \frac{\mathcal F_i}{t(x,i)}=0, \tag{2}
\]
\[
  \mathcal F_i(\mathcal P_i-1)=\mathcal D_i
  =d_x(t(x,i)-1)
  \Rightarrow
  \mathcal F_i=\frac{t(x,i)-1}{\mathcal P_i-1}d_x
  \quad (i\in G\setminus\{0\}). \tag{3}
\]
Thay phương trình (3) vào (1), (2) thu được:
\[
  \mathcal F_0+
  \left(\sum\frac{1}{\mathcal P_i-1}
  -\sum\frac{1}{t(x,i)(\mathcal P_i-1)}\right)d_x=0,
\]
\[
  \mathcal F_0=
  \left(\sum\frac{1}{t(x,i)(\mathcal P_i-1)}
  -\sum\frac{1}{\mathcal P_i-1}\right)d_x,
\]
\[
  \mathcal F_0+\sum\frac{t(x,i)-1}{\mathcal P_i-1}d_x=|G|,
\]
\[
  \left(\sum\frac{1}{t(x,i)(\mathcal P_i-1)}
  -\sum\frac{1}{\mathcal P_i-1}\right)d_x
  +\sum\frac{t(x,i)-1}{\mathcal P_i-1}d_x=|G|,
\]
\[
  \left(\sum\frac{1}{t(x,i)(\mathcal P_i-1)}
  +\sum\frac{t(x,i)}{\mathcal P_i-1}
  -2\sum\frac{1}{\mathcal P_i-1}\right)d_x=|G|.
\]
Tính trực tiếp theo công thức cuối cùng đạt $O(|G|^2)$.

Xét tối ưu. Nút thắt là ba tổng ở vế trái. Với ánh xạ
$x\mapsto 1/(\mathcal P_x-1)$, thực hiện DFT; phần thứ nhất và thứ hai lần lượt
là phần tử thứ $x$ và $x^{-1}$ của kết quả. Phần thứ ba là một hằng số. Do đó
ta tối ưu được độ phức tạp về độ phức tạp DFT, tức $O(T\log T)$, đủ để qua bài.

Đây là một ứng dụng kinh điển của DFT: dùng DFT chuyển hệ phương trình chập phức
tạp thành các phương trình nhân điểm đơn giản, rồi dùng độ dời và giá trị đặc
biệt để suy ra phương trình giải độ dời, cuối cùng giải được đáp án.

\subsection{DFT trên nửa nhóm giao hoán hữu hạn}

\subsubsection{DFT trên nhóm giao hoán hữu hạn}

\paragraph{Phân rã nhóm giao hoán hữu hạn.}

\paragraph{Định nghĩa 4.1.}
Nếu một bộ phần tử $[a_1,a_2,\ldots,a_m]$ của nhóm giao hoán hữu hạn $G$ thỏa:
\begin{itemize}
  \item $\forall x\in G,\ \exists(k_i)_{i=1}^{m}$ sao cho
  $\prod_{i=1}^{m} a_i^{k_i}=x$;
  \item $\forall(k_i)_{i=1}^{m}$, nếu $\prod_{i=1}^{m} a_i^{k_i}=e$ thì
  $a_i^{k_i}=e$ với mọi $i\in[1,m]$,
\end{itemize}
thì bộ phần tử này là một phân rã của $G$.

\paragraph{Bổ đề 4.1.}
$[a_1,a_2,\ldots,a_m]$ là một phân rã của $G$ khi và chỉ khi
\[
  (a_1)\times(a_2)\times\cdots\times(a_m)=G.
\]
Xét hai tính chất trên, chứng minh khá đơn giản.

Tiếp theo ta chứng minh mọi nhóm giao hoán hữu hạn đều tồn tại phân rã, đồng
thời đưa ra cách dựng.

\paragraph{Bổ đề 4.2.}
Mọi nhóm giao hoán hữu hạn $G$ đều có thể phân rã thành tổng trực tiếp của các
nhóm $p$ giao hoán hữu hạn.

\paragraph{Chứng minh.}
Gọi $o(x)$ là cấp của $x$. Giả sử cấp của $G$ là $p^rm$, với $p$ nguyên tố và
$p\nmid m$. Đặt
\[
  G_p=\{x\mid o(x)\mid p^r\},\qquad
  G_m=\{x\mid o(x)\mid m\}.
\]
Ta chứng minh $G_p\times G_m=G$.

Dựng ánh xạ $\varphi:(x,y)\mapsto xy$. Hiển nhiên đây là một đồng cấu nhóm; chỉ
cần chứng minh nó toàn ánh và $\ker\varphi=\{e\}$.

Với phần tử bất kỳ $x$, theo định lý Lagrange, $o(x)\mid p^rm$. Đặt
$o(x)=p^sm'$ với $p\nmid m'$, khi đó $o(x^{p^s})=m'\mid m$ và
$o(x^{m'})=p^s\mid p^r$. Vì $(p^s,m)=1$, tồn tại $u,v$ sao cho
$up^s+vm=1$, nên $x=x^{up^s+vm}=x^{up^s}x^{vm}$. Do đó $\varphi$ là toàn ánh.

Nếu $x\in G_p$, $y\in G_m$ và $xy=e$, vì
$o(x)=o(x^{-1})=o(y)$, ta có $o(x)\mid(p^r,m)=1$, nên $x=y=e$. Vậy
$\ker\varphi=\{e\}$.

Do đó $\varphi$ là một đẳng cấu nhóm. Lặp lại quá trình này, ta có thể phân rã
$G$ thành tổng trực tiếp của một số nhóm $p$ giao hoán hữu hạn. \hfill\rule{0.55em}{0.55em}

\paragraph{Phân rã nhóm $p$ giao hoán hữu hạn.}

\paragraph{Bổ đề 4.3.}
Với nhóm $p$ giao hoán hữu hạn $G$, lấy một phần tử $a$ có cấp lớn nhất và đặt
$H=(a)$. Khi đó với mọi $xH\in G/H$, tồn tại $y\in xH$ sao cho
$o(xH)=o(y)$.

\paragraph{Chứng minh.}
Đặt $o(a)=p^n$, $o(x)=p^t$, $o(xH)=p^s$; hiển nhiên $t\ge s\ge k$. Xét
$x^{p^s}=(xH)^{p^s}=H$, nên $x^{p^s}\in H$. Đặt $x^{p^s}=a^r$, khi đó
\[
  p^{t-s}=o(x^{p^s})=o(a^r)=\frac{o(a)}{(o(a),r)}
  =\frac{p^n}{(p^n,r)}.
\]
Suy ra $(o(a),r)=p^{n-t+s}$. Không mất tính tổng quát, đặt
$r=up^{n-t+s}$.

Dựng $y=xa^{p^{n-t}v}\in xH$, trong đó $v=-u p^{n-t}$. Khi đó
\[
  (yH)^{p^s}=y^{p^s}H=x^{p^s}a^{p^s p^{n-t}v}H=H.
\]
Từ đó suy ra có thể chọn đại diện trong lớp $xH$ có cấp đúng bằng $o(xH)$.
\hfill\rule{0.55em}{0.55em}

\paragraph{Bổ đề 4.4.}
Mọi nhóm $p$ giao hoán hữu hạn $G$ đều có thể được phân rã.

\paragraph{Chứng minh.}
Quy nạp theo $|G|=p^n$, giả sử mọi nhóm $G'$ có $|G'|<p^n$ đều tồn tại phân rã.
Tìm phần tử có cấp lớn nhất $a\in G$, $o(a)=p^k$. Nếu $k=n$ thì $[a]$ chính là
một phân rã.

Đặt $H=(a)$ và $G'=G/H$. Theo giả thiết quy nạp, $G'$ tồn tại phân rã; giả sử
là $[a_1H,a_2H,\ldots,a_mH]$. Theo Bổ đề 4.3, chọn các đại diện sao cho
$o(a_i)=o(a_iH)$ với $i\in[1,m]$.

Ta chứng minh $[a,a_1,a_2,\ldots,a_m]$ là một phân rã.

Với mọi $x\in G$, xét $xH\in G/H$. Có
\[
  xH=\prod_{i=1}^{m}(a_iH)^{k_i}
  =\left(\prod_{i=1}^{m}a_i^{k_i}\right)H.
\]
Do đó tồn tại $t$ sao cho $x=a^t\prod_i a_i^{k_i}$.

Ngược lại, nếu $a^t\prod_i a_i^{k_i}=e$, thì
$\prod_i(a_iH)^{k_i}=H$. Theo tính chất của $a_iH$, suy ra
$(a_iH)^{k_i}=H$ với mọi $i$, tức $a_i^{k_i}\in H$. Vì
$o(a_i)=o(a_iH)$, ta có $a_i^{k_i}=e$. Khi đó cũng suy ra $a^t=e$.
\hfill\rule{0.55em}{0.55em}

\paragraph{Định lý 4.1.}
Mọi nhóm giao hoán hữu hạn đều có thể được phân rã.

Điều này suy ra ngay từ Bổ đề 4.2 và Bổ đề 4.4.

\paragraph{Một cách cài đặt đơn giản hơn.}

Chứng minh trên thực ra đã cho một quá trình dựng, nhưng vẫn quá phức tạp.

\paragraph{Định lý 4.2.}
Với nhóm giao hoán hữu hạn $G$, duy trì dãy $B$, ban đầu rỗng, và
$S=\operatorname{span}(B)$. Đặt $f(x)=\min\{k\mid x^k\in S\}$. Lặp các thao tác:
\begin{itemize}
  \item nếu $G=S$ thì kết thúc;
  \item tìm mọi $x\in G$, $x\notin S$, và chọn phần tử $a$ có $o(x)/f(x)$ lớn
  nhất;
  \item $B:=B+a$, $S=\operatorname{span}(B)$.
\end{itemize}
Dãy $B$ thu được là một phân rã.

Xét mệnh đề mạnh hơn sau.

\paragraph{Bổ đề 4.5.}
Với mọi $n\ge0$, tại một thời điểm nào đó nếu
$B=[a_1,a_2,\ldots,a_n]$, thì tồn tại một phân rã
$[a_1,a_2,\ldots,a_m]$ ($m\ge n$) sao cho
$o(a_i)=p_i^{k_i}$ với $n<i\le m$, $p_i$ nguyên tố và $k_i\in\mathbb Z^+$.

\paragraph{Chứng minh.}
Quy nạp theo $n$; với $n=0$ hiển nhiên đúng. Giả sử trước khi phần tử thứ $n$
được thêm vào, một phân rã hợp lệ là
$B'=[a_1,\ldots,a'_r]$.

Với phần tử $x$, nếu $x=\prod_i (a'_i)^{t_i}$, thì
\[
  f(x)=\operatorname{lcm}_i\frac{o(a'_i)}{(o(a'_i),t_i)}.
\]
Khi $f(x)$ lớn nhất, có thể tìm một lũy thừa $p^k$ và một chỉ số $i$ sao cho
$(o(a'_i)/p^k,t_i)=1$ và $k$ đạt lớn nhất; từ đó dựng được một phần tử $x$ có
$o(x)=f(x)$, $x\notin S$ và $f(x)=\max_{y\in G\setminus S} f(y)$.

Dựng
\[
  B=(a_1,\ldots,a_{n-1},x)+R,
\]
trong đó $R$ nhận được bằng cách thay trong phân rã cũ một phần tử tương ứng bởi
các phần còn lại. Không khó chứng minh $\operatorname{span}(B)$ chứa phần tử bị
xóa, nên $\operatorname{span}(B)=G$; hơn nữa cấp của các phần tử trong $B$ vẫn
đúng dạng trên, vì vậy đây là một phân rã. \hfill\rule{0.55em}{0.55em}

\paragraph{Bổ đề 4.6.}
Với cài đặt thích hợp, cách làm này có độ phức tạp $O(|G|^2)$.

\paragraph{Chứng minh.}
Nút thắt độ phức tạp nằm ở việc tính $f(x)$. Ở vòng thứ $n$, ta có
\[
  f(x)\le\prod_i o(a_i),\qquad
  f(x)<\frac{|G|}{\prod_i o(a_i)}.
\]
Vì vậy tổng chi phí tính $f(x)$ không vượt quá $O(|G|)$ cho mỗi phần tử, dẫn
tới tổng $O(|G|^2)$. \hfill\rule{0.55em}{0.55em}

\paragraph{Nửa nhóm giao hoán hữu hạn.}

\paragraph{Định lý 4.3.}
$G\to\mathbb C$ tồn tại DFT khi và chỉ khi $G$ thỏa luật tuần hoàn.

Chứng minh và cách dựng được trình bày trong tài liệu~[6], bài viết này không
triển khai chi tiết.

\subsection{Nhóm đối xứng}

\subsubsection{Định nghĩa}

\paragraph{Định nghĩa 5.1.}
Định nghĩa tích của hai biến đổi tuyến tính là tác động liên tiếp của hai biến
đổi tuyến tính, tức phép nhân ma trận. Khi đó, với không gian véc-tơ $V$ $n$
chiều, mọi biến đổi tuyến tính không suy biến, tức có định thức khác $0$, tạo
thành nhóm tuyến tính tổng quát $n$ chiều, ký hiệu là $GL(V)$. Nhóm con $L(V)$
của nó được gọi là nhóm biến đổi tuyến tính trên $V$.

\paragraph{Định nghĩa 5.2.}
Với nhóm $G$, nếu tồn tại một đồng cấu $p$ từ $G$ tới nhóm biến đổi tuyến tính
$L(V)$ trên không gian tuyến tính $V$, thì $p$ được gọi là một biểu diễn tuyến
tính của nhóm $G$. Không gian $V$ là không gian biểu diễn; số chiều của $V$ là
số chiều của biểu diễn, ký hiệu là $d_p$.

\paragraph{Định nghĩa 5.3.}
Giả sử hai biểu diễn $p,p'$ của nhóm $G$ trên không gian $V$ lần lượt lấy cơ sở
\[
  e=(e_1,e_2,\ldots,e_n),\qquad
  e'=(e'_1,e'_2,\ldots,e'_n).
\]
Nếu $e=e'X$, $\det X\ne0$, và với mọi $g\in G$ có
$p(g)=X^{-1}p'(g)X$, thì gọi $p,p'$ là tương đương.

\paragraph{Định nghĩa 5.4.}
Nếu $p$ là một biểu diễn của nhóm $G$ trên không gian tuyến tính $V$, và không
tồn tại không gian con $W\subset V$ thỏa $W\ne0$, $W\ne V$ và với mọi
$g\in G$, $x\in W$ đều có $p(g)x\in W$, thì $p$ là một biểu diễn bất khả quy.

Ký hiệu tập mọi biểu diễn bất khả quy đôi một không tương đương của $G$ là
$R(G)$.

\subsubsection{DFT của nhóm đối xứng}

Tiếp theo bài viết đưa ra thuật toán DFT cho nhóm $S_n$. Do giới hạn dung lượng
và trình độ của tác giả, phần này không đưa ra chứng minh và cài đặt cụ thể;
xem các tài liệu~[1], [2], [3].

\paragraph{Định lý 5.1 (định lý Burnside).}
Với nhóm hữu hạn $G$, ta có
\[
  \sum_{p\in R(G)} d_p^2=|G|.
\]
Định lý này chứng minh rằng sau DFT, lượng thông tin là đủ.

\paragraph{Bổ đề 5.1.}
Với nhóm hữu hạn $G$, cách dựng sau là một DFT hợp lệ:
\begin{itemize}
  \item $a_i=d_{p_i}$;
  \item $\mathcal F(A)_i=\sum_{g\in G} A_g p_i(g)$.
\end{itemize}
Điều này cho thấy biểu diễn nhóm và DFT có liên hệ rất chặt. Nếu đã biết
$R(G)$, ta có thể hoàn thành DFT trong $O(|G|^2)$. Đáng tiếc là hiện không có
một thuật toán tổng quát để tìm $R(G)$ của nhóm hữu hạn tùy ý $G$.

\paragraph{Bổ đề 5.2.}
Với cách dựng trên, một cách tính IDFT khác là
\[
  A_x=\frac{1}{|G|}\sum_{p_i\in R(G)}
  d_{p_i}\operatorname{Tr}\!\left(p_i(x^{-1})\mathcal F(A)_i\right).
\]

\paragraph{Bổ đề 5.3.}
$R(S_n)$ có quan hệ một-một với các phân hoạch của $n$.

Tiếp theo ta đưa ra một cách dựng gọi là biểu diễn bán chuẩn Young.

\paragraph{Định nghĩa 5.5.}
Với hai bảng Young chuẩn bậc $n$ có cùng hình dạng $T_1,T_2$, định nghĩa thứ tự
không nhãn của chúng như sau. Nếu số $n$ nằm ở hai hàng khác nhau trong
$T_1,T_2$, thì $T_1<T_2$ khi và chỉ khi hàng của $n$ trong $T_1$ nhỏ hơn. Nếu
không, đặt $T'_1,T'_2$ là hai bảng Young thu được sau khi xóa phần tử $n$ khỏi
$T_1,T_2$, và đặt $T_1<T_2$ khi và chỉ khi $T'_1<T'_2$.

\paragraph{Định nghĩa 5.6.}
Với một bảng Young $T$, giả sử phần tử $n,m$ lần lượt nằm ở các hàng
$x_1,x_2$ và các cột $y_1,y_2$. Định nghĩa khoảng cách trục
\[
  d(n,m)=x_2-x_1+y_1-y_2,
\]
chú ý không lấy trị tuyệt đối.

\paragraph{Bổ đề 5.4.}
Hình dạng của các bảng Young bậc $n$ tương ứng một-một với các phân hoạch của
$n$. Chỉ cần xét độ dài từng hàng của bảng Young là thấy.

\paragraph{Định nghĩa 5.7.}
Với phân hoạch $a$ của $n$, sắp các bảng Young tương ứng theo thứ tự không nhãn
thành $T$. Ký hiệu khoảng cách trục của $a,b$ trong $T_i$ là $d_i(a,b)$. Dưới
đây cho cách dựng $p((t,t+1))$; đặt $p((t,t+1))=\alpha$.
\begin{itemize}
  \item Với mọi $i$, $\alpha_{i,i}=d_i(t,t+1)^{-1}$.
  \item Nếu đổi chỗ $t$ và $t+1$ trong $T_i$ vẫn là bảng Young chuẩn hợp lệ, gọi
  bảng thu được là $T_j$, thì
  \[
    \alpha_{i,j}=
    \begin{cases}
      1-d_i(t,t+1)^{-2},& i<j,\\
      1,& i>j.
    \end{cases}
  \]
  \item Các vị trí còn lại đều đặt bằng $0$.
\end{itemize}
Với hoán vị bất kỳ $p$, dùng sắp xếp nổi bọt có thể chứng minh nó phân rã được
thành tích của một số hoán vị dạng $(i,i+1)$; từ đó hợp nhất để thu được
$p(p)$. Có thể chứng minh đây là một họ biểu diễn bất khả quy và đôi một không
tương đương. Tính trực tiếp, chi phí DFT là $O(|G|^2)$.

\subsubsection{Độ phức tạp tốt hơn}

Xét $S_n$ và đặt
\[
  H=\{x\mid x(n)=n\}\sim S_{n-1}.
\]
Khi đó $\{s_iH\mid s_i\in S_n\}$ tạo thành một phân hoạch theo các lớp kề của
$S_n$. Biến đổi công thức định nghĩa DFT:
\[
\begin{aligned}
\mathcal F(A)
&=\sum_{g\in G} A_gp_i(g)\\
&=\sum_t p_i(s_t)\sum_{x\in H} A_{s_tx}p_i(x).
\end{aligned}
\]
Đặt $A'_x=A_{s_tx}$.

\paragraph{Định lý 5.2.}
Với $x\in H$, dạng của $p(x)$ là
\[
\begin{bmatrix}
p_{i_1}(x)&0&\cdots&0\\
0&p_{i_2}(x)&\cdots&0\\
0&0&\ddots&p_{i_k}(x)
\end{bmatrix},
\]
trong đó $i_j$ biểu thị hình dạng phân hoạch hoặc bảng Young bậc $n-1$ có thể
thu được sau khi xóa một phần tử khỏi hình dạng phân hoạch hoặc bảng Young
$a_j$.

Chỉ cần khảo sát vị trí của $n$ trong bảng Young. Các bảng có cùng vị trí của
$n$ chắc chắn liên tiếp trong $T$. Vì $x(n)=n$, xóa $n$ không ảnh hưởng quá
trình dựng. Quy nạp là chứng minh được.

Do đó với $A':S_n\to\mathbb C$, ta có thể giải riêng từng $\mathcal F(A')$, rồi
đệ quy thành bài toán con $S_{n-1}\to\mathbb C$, sau đó hợp nhất bằng vài lần
nhân dãy ma trận.

\paragraph{Bổ đề 5.5.}
Trong các ma trận
$C_t=p((t,t+1))$ với $t\in[1,n-2]$ và
$D_t=D_{\mathrm{prev}}A_t p(t)$ với $t\in[1,n)$, tổng số phần tử khác $0$ là
$O(n!)$.

\paragraph{Chứng minh.}
Xét mỗi phân hoạch hoặc hình dạng bảng Young $a$ của $n-1$. Nó đóng góp nhiều
nhất một lần vào mỗi $C_t$ và $D_t$, tức xuất hiện nhiều nhất $2n-2$ lần.
Ta biết $\sum a_i=(n-1)!$, nên tổng đóng góp không vượt quá $O(n!)$.
\hfill\rule{0.55em}{0.55em}

Với $s_i$, ta chọn $s_i=(n,n-1,\ldots,i)$. Khi đó $s_i=s_{i-1}\times(i,i+1)$;
như vậy chỉ cần $n$ lần nhân dãy ma trận là tính được $p(s_i)$. Chú ý rằng nhờ
Bổ đề 5.5, tổng độ phức tạp của phép cộng và phép nhân ma trận không vượt quá
$O(n!\,n)$.

Vì vậy độ phức tạp DFT của $S_n$ là
\[
  T(n)=nT(n-1)+O(n!\,n^2)=O(n!\,n^3),
\]
tốt hơn rõ rệt so với cách thô bạo $O((n!)^2)$.

Thuật toán này dùng chia để trị tương tự FFT, đệ quy các trường hợp khác nhau
về cùng một bài toán con để giảm độ phức tạp. Với các nhóm tổng quát hơn, nếu
tìm được cách dựng tương tự, ta cũng có thể hoàn thành DFT với độ phức tạp thấp.

\subsection*{Kết luận}

DFT đã được dùng rộng rãi trong OI, nhưng trong thời gian dài ứng dụng của DFT
hầu như chỉ giới hạn ở phép chập dãy; những năm gần đây mới mở rộng tới chuỗi
lũy thừa tập hợp. Tuy nhiên với các định nghĩa phép chập tổng quát hơn, cộng
đồng OI chưa nghiên cứu nhiều. Bài viết này cho thấy nhóm không giao hoán cũng
có thể thực hiện phép chập với độ phức tạp cao hơn, đồng thời chỉ ra quan hệ
giữa DFT và biểu diễn nhóm. Tác giả chỉ trình bày DFT của nhóm đối xứng, chưa
trình bày ứng dụng của nó; ý nghĩa sâu hơn và ứng dụng đa dạng hơn của DFT vẫn
cần các bạn đọc quan tâm tiếp tục nghiên cứu.

\subsection*{Lời cảm ơn}

Tác giả cảm ơn China Computer Federation đã cung cấp nền tảng học tập và trao
đổi; cảm ơn thầy Xu Xianyou đã quan tâm và hướng dẫn; cảm ơn bạn Xu Tingqiang
và anh Ma Yaohua đã giúp đỡ về nội dung; cảm ơn bạn Xu Tingqiang, bạn Luo Kai,
anh Jin Ce và anh Xu Yinzhan đã đọc kiểm tra bản thảo.

\subsection*{Tài liệu tham khảo}

\begin{enumerate}
  \item Michael Clausen; Ulrich Baum. \emph{Fast Fourier transforms for
  symmetric groups: theory and implementation}. Math. Comp. 61 (1993),
  pp. 833--847, 1993.
  \item Risi Kondor. \emph{A C++ library for fast Fourier transforms on the
  symmetric group}.
  \item Persi Diaconis; Daniel Rockmore. \emph{Efficient computation of the
  Fourier transform on finite groups}. J. Amer. Math. Soc. 3 (1990),
  pp. 297--332, 1990.
  \item Audrey Terras. \emph{Fourier Analysis on Finite Groups and
  Applications}. Cambridge University Press, 1999.
  \item Josh Alman; Virginia Vassilevska Williams. \emph{A refined laser method
  and faster matrix multiplication}. The 32nd Annual ACM-SIAM Symposium on
  Discrete Algorithms, 2020.
  \item Liu Cheng'ao. \emph{Bàn về ứng dụng của DFT trong thi tin học}. Tập
  luận văn đội tuyển ứng viên quốc gia Trung Quốc IOI 2018, 2018.
\end{enumerate}

\section{Chủ đề chính}

Từ mục lục đã đọc được trên ảnh render, tập năm 2022 bao phủ các nhóm chủ đề
sau:
\begin{itemize}
  \item duy trì thông tin đường đi trên DAG, DFT và các biến thể tổng quát;
  \item tư tưởng tham lam dựa trên so sánh và thuật toán xấp xỉ;
  \item đồ thị phẳng, tô màu đồ thị và các bài toán tương tác dạng tìm tham số;
  \item tối ưu không gian, ma trận Monge và phân tích bảng băm;
  \item đại số tuyến tính, lý thuyết Lyndon, tổ hợp chuỗi và mô phỏng luồng chi
  phí.
\end{itemize}

\section{Ghi chú về trích xuất}

\texttt{pdftotext} không trích xuất được văn bản có nghĩa từ tập 2022; kết quả
kiểm tra các trang đầu chỉ gồm dấu sang trang. Mục lục và bài đầu tiên trong bản
dịch này được đọc bằng cách render bằng \texttt{pdftoppm} rồi OCR bằng
\texttt{tesseract} với ngôn ngữ \texttt{chi\_sim+eng}. Một số công thức ở bài
đầu tiên, cụ thể các mục 2.1, 2.3, 3.2 và 3.3, bị nhiễu trên ảnh/OCR; các chỗ
đó được đánh dấu bằng ghi chú OCR tiếng Việt thay vì tự suy diễn công thức.
Bài thứ hai được OCR từ các trang vật lý 17--31 của PDF, tương ứng trang mục
lục 13--27; các khối công thức quan trọng được đối chiếu lại trên ảnh render.

\end{document}
